% basic_principles_symbols.tex -- Book I. Symbols and Abbreviations (jemmybutton-matched).
% Text from jemmybutton byrne-en-latex.tex (CC-BY-SA).
% Build: lualatex basic_principles_symbols.tex  ->  basic_principles_symbols.pdf
\documentclass[12pt]{article}
\usepackage{geometry}
\geometry{a5paper, margin=1.5cm}
\usepackage[lining]{ebgaramond}
\usepackage{amsmath}
\usepackage{amssymb}
\usepackage{byrne}
\def\mpPre{u := 1cm; textLabels := false;}

% Symbol row: left=symbol (22%), right=description (74%)
\newcommand{\symrow}[2]{%
  \vspace{0.25cm}%
  \noindent%
  \begin{minipage}[t]{0.22\linewidth}%
    \vspace{0pt}\centering\Large #1%
  \end{minipage}\hfill%
  \begin{minipage}[t]{0.74\linewidth}%
    \vspace{4pt}#2%
  \end{minipage}%
}

\begin{document}

\begin{center}
  {\large\textbf{Book~I.}}\\[4pt]
  {\Large\textbf{Symbols and Abbreviations}}
\end{center}

\vspace{0.4cm}
\noindent\rule{\linewidth}{0.4pt}
\vspace{0.2cm}

\symrow{$\therefore$}{expresses the word \emph{therefore}.}

\symrow{$\because$}{expresses the word \emph{because}.}

\symrow{$=$}{expresses the word \emph{equal}. This sign of equality may be
  read \emph{equal to}, or \emph{is equal to}, or \emph{are equal to}; but
  the discrepancy in regard to the introduction of the auxiliary verbs
  \emph{is}, \emph{are}, \&c.\ cannot affect the geometrical rigour.}

\symrow{$\neq$}{means the same as if the words \emph{`not equal'} were written.}

\symrow{$>$}{signifies \emph{greater than}.}

\symrow{$<$}{signifies \emph{less than}.}

\symrow{$\ngtr$}{signifies \emph{not greater than}.}

\symrow{$\nless$}{signifies \emph{not less than}.}

\symrow{$+$}{is read \emph{plus} (\emph{more}), the sign of addition; when
  interposed between two or more magnitudes, signifies their sum.}

\symrow{$-$}{is read \emph{minus} (\emph{less}), signifies subtraction; and
  when placed between two quantities denotes that the latter is taken from
  the former.}

\symrow{$\times$}{expresses the product of two or more numbers when placed
  between them in arithmetic and algebra; but in geometry it is generally
  used to express a \emph{rectangle}, when placed between ``two straight
  lines which contain one of its right angles.'' A \emph{rectangle} may also
  be represented by placing a point between two of its conterminous sides.}

\symrow{$: :: :$}{expresses an \emph{analogy} or \emph{proportion}; thus if
  A, B, C and D represent four magnitudes, and A has to B the same ratio
  that C has to D, the proportion is written
  $A : B :: C : D$,\ \ $A : B = C : D$,\ \ or\ \ $\dfrac{A}{B} = \dfrac{C}{D}$.\\[4pt]
  This equality is read: \emph{as A is to B, so is C to D}.}

\symrow{$\parallel$}{signifies \emph{parallel to}.}

\symrow{$\perp$}{signifies \emph{perpendicular to}.}

% --- Angle and right-angle symbols (need a figure) ---
\defineNewPicture{%
  pair A, B, C, D;%
  numeric s;%
  s := 3/2u;%
  A := (0, 0);%
  B := dir(0)*s;%
  C := dir(50)*s;%
  D := dir(90)*s;%
  byAngleDefine(B, A, C, byblack, ARC_SECTOR);%
  byAngleDefine(B, A, D, byblack, ARC_SECTOR);%
  byPointLabelRemove(A,B,C,D);%
}

\symrow{\drawAngle{BAC}}{signifies \emph{angle}.}

\symrow{\drawAngle{BAD}}{signifies \emph{right angle}.}

\symrow{\normalsize\drawTwoRightAngles}{signifies \emph{two right angles}.}

% --- Point symbols (need a second figure) ---
\defineNewPicture{%
  pair A, B, C, D;%
  A := (0, -1/4u);%
  B := (u, 0);%
  C := (-u, 0);%
  D := (0, u);%
  byLineDefine(A, D, byblack, SOLID_LINE, REGULAR_WIDTH);%
  byLineDefine(B, D, byblack, SOLID_LINE, REGULAR_WIDTH);%
  byLineDefine(C, D, byblack, SOLID_LINE, REGULAR_WIDTH);%
  byPointLabelRemove(A,D);%
}

\symrow{%
  \drawFromCurrentPicture{%
    draw byNamedLine(AD);%
    draw byNamedLineSeq(0)(BD,CD);%
  }
  \quad or \quad
  \drawFromCurrentPicture{%
    draw byNamedLineSeq(0)(AD,BD);%
  }%
}{briefly designates a \emph{point}.}

\vspace{0.4cm}

\noindent
The square described on a line is concisely written thus,
$\drawUnitLine{AD}^2$.
In the same manner twice the square of, is expressed by
$2 \cdot \drawUnitLine{AD}^2$.

\vspace{0.5cm}
\noindent\rule{\linewidth}{0.4pt}
\vspace{0.3cm}

\symrow{def.}{signifies \emph{definition}.}

\symrow{post.}{signifies \emph{postulate}.}

\symrow{ax.}{signifies \emph{axiom}.}

\symrow{hyp.}{signifies \emph{hypothesis}. The \emph{hypothesis} is the
  condition assumed or taken for granted. Thus, the hypothesis of the
  proposition that a triangle is isosceles is that its legs are equal.}

\symrow{const.}{signifies \emph{construction}. The \emph{construction} is
  the change made in the original figure, by drawing lines, making angles,
  describing circles, \&c., in order to adapt it to the argument of the
  demonstration or the solution of the problem. The conditions under which
  these changes are made are as indisputable as those contained in the
  hypothesis.}

\symrow{Q.~E.~D.}{signifies \emph{Quod erat demonstrandum} — which was to
  be demonstrated.}

\end{document}
